% No preamble. Prefix qll-. Actual homogeneous local sector, fixed classical system.
\section{A local limit and quotient descent of the quantum macrostep}
\label{qll-section}

Fix the composite ratio \(R\), auxiliary grouping \(J\), preceding
homogeneous depth \(M\), real retained link/color windows, cutoff
threshold \(T\), positive coupling \(g\), and common field radius
\(b\) from the common-domain construction. All fluctuation-cell
dimensions are fixed once \(R,J\) are fixed. Compatible cubic
tori exhaust the current and next lattices together. Write
\[
 \|E\|_{a_0,b}
   =\sup_i\sum_{H\ni i}e^{a_0|H|}\|E_H\|_b,
 \qquad a_0=1/16,
\]
for the connected activity norm. Quotient classes identify
vacuum-normalized totals on compact fields, as in Notebook29.

The actual finite local correction has the decomposition
\begin{align}
 \mathcal F_N(E_N)
   &=\mathcal T_NE_N+\Delta_N+\mathcal Q_{{\rm det},N}
                         +\mathcal L_N(Q_{E_N}),\notag\\
 \mathcal G_N(E_N)
   &=\Delta_N+\mathcal Q_{{\rm det},N}
                         +\mathcal L_N(Q_{E_N}).
 \label{qll-finite-map}
\end{align}
Here \(\mathcal T_N\) is pre-gauge background substitution,
\(Q_{E_N}\) is the inherited fluctuation difference in the
actual chart, and \(\mathcal L_N\) is the strict Gaussian
cluster logarithm with its fixed classical base factors and
precision. The pivot term, determinant, sharp windows and
field-independent normalization are those of the finite
common-domain integral. Only after constructing this expression
do we subtract its value at zero background to form a quotient
total. A separately tracked vacuum normalization is not
identified with zero.

\subsection{One fixed, input-independent classical system}
\label{qll-data-section}

Call coefficients \emph{locally normal} when every fixed lifted
finite carrier coefficient converges uniformly on compact
subsets of its finite-dimensional background polydisc.
For joint coefficients require this convergence uniformly on
the fixed real window of their finitely many fluctuation
coordinates. Matrix coefficients use their finite block norm.
The final norm bound is on the original open \(b\) ball;
local normal convergence is only asserted on compact subsets.

Choose one common regulator subsequence with the following
actual classical data.

\begin{enumerate}
\item Every fixed baseline covariance entry \(C_{0,N}(i,j)\)
converges. The baseline has the uniform exponential Schur
bounds and positive spectral lower bound used in the strict
Gaussian theorem.
\item The strict precision coefficients \(B_{N,Y}\), classical
base exponents, pivot-log factors, and coherent pre-gauge
and joint-chart support coefficients converge locally normally.
Both matrix indices of \(B_{N,Y}\) belong to \(Y\).
\item All stronger support bounds, analytic neighborhoods,
strict sufficient budgets and exact local scalar identities
are those of the established finite construction. Aggregate
primitive labels by carrier and type before taking limits:
there are only finitely many types per fixed carrier, including
precision, base, inherited and bad-block types.
\end{enumerate}

These data are chosen independently of the inherited interaction.
At fixed parameters their strict families admit a countable
diagonal extraction. Each fixed coefficient has finitely many
coordinates. For the real fluctuation window, use the actual
unwindowed chart's holomorphic extension to a common complex
neighborhood and its uniform Cauchy bounds. Bounded measurable
functions alone would not give this compactness. The sharp
indicators remain fixed; they are not differentiated or subjected
to a complex change of variables.

The actual flat formulas already give compatible periodized
baseline kernels. More generally, joint entrywise extraction
with the stated uniform bounds preserves every finite
principal spectral inequality. The conclusions below refer to
the selected limiting classical system, not to an arbitrary
different sequence of baselines or classical branches.

\subsection{Explicit strict transport margins}
\label{qll-margins-section}

The common-domain proof permits the following stricter choices.
Fix its \(B=53/8\), primitive constants and classical support
reserves first. Choose \(R\) so that
\begin{equation}
 c_1(R)B\le a_0/4,
 \qquad c_1(R)=\frac{5^d2^d}{k_R},
 \qquad k_R=1+\lfloor(R-3)/4\rfloor .
 \label{qll-R-margin}
\end{equation}
After fixing \(R\), choose \(J\) and its logarithmic tail
collar so that, with the actual tail-to-margin ratio \(t\),
\begin{equation}
 t<1,\qquad
 D(J)\log\frac1{1-t}\le a_0/4.
 \label{qll-J-margin}
\end{equation}
This is possible because its sufficient bound is a constant
times \(J^{2d}e^{-\kappa_0\ell_{\rm tail}/2}\), which tends
to zero with the established choice of collar. It is independent
of the later choice of \(b\). Thus
\begin{equation}
 \delta_{\rm card}
   :=a_0-c_1(R)B-D(J)\log\frac1{1-t}
   \ge a_0/2>0.
 \label{qll-card-margin}
\end{equation}
The original weak half allocations alone would only give a
nonnegative margin; \eqref{qll-R-margin}--\eqref{qll-J-margin}
make the needed strictness explicit.

In the positive joint core/tail majorant, restricting original
carriers to \(|H|>Q\) now gives the bound
\begin{equation}
 C\epsilon e^{-\delta_{\rm card}Q}.
 \label{qll-input-carrier-tail}
\end{equation}
Choose also a fixed \(\zeta>1\) sufficiently close to one that
\begin{equation}
 \zeta t<1,\qquad
 c_1(R)B+D(J)\log\frac1{1-\zeta t}<a_0.
 \label{qll-tail-degree-margin}
\end{equation}
Replacing \(t\) by \(\zeta t\) in the exact scalar-multiindex
majorant gives an \(O(\zeta^{-k})\) bound above tail-expansion
degree \(k\). If \(t=0\), there are no positive tail degrees.
The rooted spatial weight and the strict spatial budget
\(\kappa_{\rm tail}>B+\delta\) control omitted remote supports.
The whole output-fiber maximum remains inside the support sum.

Choose \(b\), then the permitted \(T,g\), within the common-domain
order with strict slack in the precision, density, bad-block
and multiplier inequalities. Equivalently use the smaller
quarter budgets of the two-input comparison. Strictness allows
a fixed bookkeeping dilation greater than one in the positive
majorants below. All these choices are uniform in volume and
in preceding \(M\) when the established classical estimates
are uniform there.

\subsection{Continuity of the actual strict Gaussian construction}
\label{qll-Gaussian-section}

\begin{lemma}[Local construction continuity]
\label{qll-Gaussian-lemma}
Use fixed finite-dimensional fluctuation cells, fixed real
windows, the classical data above, and locally normal multiplier
coefficients under uniform strict budgets. Fix deterministic
connecting graphs that agree on every bounded lifted
configuration for sufficiently large periods. The canonical
connected Gaussian logarithm coefficients then have unique
locally normal limits for those limiting local data and
assignments. Their original anchored bound holds on the full
open \(b\) ball.
\end{lemma}

\begin{proof}
The proof retains the construction's positive majorants, not
merely its final activity norm.

First expand the precision inverse in its \(B_Y\) labels.
A term has the form
\[
 C_0(b,u_1)B_{Y_1}(u_1,v_1)C_0(v_1,u_2)
       \cdots B_{Y_r}(u_r,v_r)C_0(v_r,c).
\]
Every intermediate index belongs to a retained precision
carrier, because both indices of \(B_Y\) lie in \(Y\).
The endpoints \(b,c\) belong to the finite moment-root set
\(\mathcal M\). If \(X=\bigcup_jY_j\), the final connected
carrier contains \(\mathcal M\cup X\). Masked baseline
inverses act on the finite set \(\mathcal M\). External
determinant trace chains have the same property. Consequently
a fixed final carrier bounds every primitive spatial label
and covariance-chain index. No hidden unbounded index excursion
or new spatial reserve is needed. The full \(C_0\) kernel entry
is used as an entry, not as an unrecorded walk expansion.

Next keep tags on primitive positive weights and dilate them:
\begin{equation}
 B_Y\mapsto zB_Y,\qquad
 w_Y\mapsto zw_Y,\qquad q_i\mapsto zq_i,
 \qquad |z|\le1+\delta_0
 \label{qll-mark-dilation}
\end{equation}
for a sufficiently small fixed \(\delta_0>0\).
Here \(w_Y\) are interaction multipliers and \(q_i\) are
bad-block factors. Strict budgets make the same positive
overlap, endpoint-tree and Mayer majorants summable on this
disc. Scaling a bad-block factor only scales its joint-tail
activity weight; no new cutoff or differentiability assertion
is involved.

Every nonempty logarithm polymer contains an interaction or
bad-block mark. Mayer repetitions repeat such marks. Covariance
forest edges number at most the molecule count minus one,
while each inverse or determinant density order counts its
precision occurrences. Thus the remaining combinatorial tails
are \(O((1+\delta_0)^{-k})\), together with the strict geometric
precision-order tail if counted separately. There is no
unmarked unbounded combinatorial order left.

For a fixed final carrier and bounded orders there are only
finitely many discrete choices, since primitive labels have
already been aggregated by carrier and type. A retained term
depends on finitely many baseline entries and analytic
precision or multiplier coefficients. Its Gaussian moment
uses a finite principal covariance. The partition masks and
degree-adapted endpoint perturbations remain in the common
accretive neighborhood. Matrix determinants and inverses are
continuous holomorphic functions there. The established
density majorant is integrable on the entire fixed real
Gaussian space, uniformly in the regulator.

Dominated convergence therefore passes the Gaussian integrals
and compact forest-parameter integrals. The fixed sharp
indicators are only bounded measurable factors; they remain
untouched. Multiplier convergence is uniform on their compact
retained windows. The argument is uniform on compact background
polydiscs, or equivalently follows there from the common
holomorphic bounds and uniqueness of the limit.

The order-tail bounds now pass each complete coefficient.
Deterministic hull grouping preserves this convergence.
Winding configurations cannot remain in a fixed bounded
lifted carrier as the period grows. The final anchored
cardinality bound controls the omitted large carriers when
passing normalized totals on compact fields.

Finally apply Fatou to finite carrier selections at radius
\(r<b\), then increase \(r\) to \(b\) and exhaust the carriers.
This retains the original full open-domain bound. Compact-total
convergence follows by the same cardinality-tail argument as
analytic interaction compactness. The result is uniqueness of
the local limit for the selected local data and assignments,
not convergence in the global anchored norm.
\end{proof}

\subsection{Inherited insertions and external factors}
\label{qll-inherited-section}

Let finite input families converge locally normally with a
uniform activity budget. The core/tail construction evaluates
each input coefficient at an analytic core in a fixed strict
interior of its old \(b\) polydisc, then expands only in its
small supported tail. It retains the projected input carrier
and every attached background or fluctuation support.

Truncate original carrier size, remote support extent and
tail-expansion degree. Equations
\eqref{qll-input-carrier-tail}--\eqref{qll-tail-degree-margin}
and the rooted exponential bound give uniform positive tails.
On the remaining finite list, ordinary Cauchy estimates pass
all needed input derivatives and evaluations. The actual core
and fluctuation chart remain in their fixed interior margin,
uniformly on the real retained windows. The complete-fiber
norm prevents an extra output-coordinate count. Hence
\(Q_{E_N}\) has locally normal coefficient limits. The
compatible assignment is linear in the inherited input.
Equal-carrier exponentiation then preserves local convergence
and the required multiplier bounds.

Lemma~\ref{qll-Gaussian-lemma} applies to \(\mathcal L_N\).
The external pivot factors pass by local analytic composition.
The external determinant term passes by its absolutely
convergent closed-chain trace series, with every precision
label and its indices retained. Classical substitution passes
by the proved summable-field composition limit. We conclude
that \eqref{qll-finite-map} has locally normal connected
coefficient limits under these hypotheses.

\subsection{Canonical periodization and quantum quotient descent}
\label{qll-descent-section}

The periodization theorem supplies a covariant representative
of every translation-invariant class without norm increase,
a connected finite-periodic representation at the same
\(a_0,b\), and exact independence of that representative as
a finite vacuum-normalized total. It also supplies finite
local gauge and spatial identities on the common chart.
For a fixed covariant representative, its fixed lifted
coefficients stabilize exactly by
\eqref{pid-fixed-carrier}; in particular they converge locally
normally. On covariant representatives the orbit construction
is linear, including on differences.

\begin{theorem}[Quantum map on the selected local quotient]
\label{qll-descent-theorem}
Under the fixed classical-system, strict-margin and domain
hypotheses above, canonical periodized inputs define a
well-defined limiting map \(\mathcal F_g\) on a sufficiently
small ball of invariant analytic interaction classes.
Its output has a connected representation at the same
\(a_0,b\). The map preserves the local total gauge and spatial
identities of the actual sector. Writing \(\mathcal T\) for
classical substitution, the quantum remainder is
\(\mathcal G_g=\mathcal F_g-\mathcal T\).
\end{theorem}

\begin{proof}
Choose an allowed covariant representative \(E\). Its finite
periodizations and the fixed classical data meet the preceding
local-continuity hypotheses. They therefore produce a limiting
connected family \(\mathcal F(E)\).

If \(E'\) represents the same class, its finite periodized
normalized total is exactly the same. Their actual inherited
fluctuation differences agree as scalar insertions, and their
background terms agree after the chosen normalization.
If raw conventions differ, their constant scalar factor is
handled in the separate vacuum normalization. The fixed
Gaussian, windows and chart are otherwise identical.
Consequently their actual finite output normalized totals
agree, with compatible nonvanishing logarithm branches.
The canonical coefficient arrays need not agree.

Passing their local limits and then compact normalized totals
shows \([\mathcal F(E)]=[\mathcal F(E')]\). This proves descent.
The exact finite local gauge and spatial identities pass on
compact backgrounds and their allowed transforms. Holomorphic
continuation gives their local complex identities on the
common ball. Norm bounds pass by Fatou. The same reasoning
applies to \(\mathcal T\) and hence to the indicated remainder.
\end{proof}

\subsection{A quotient Lipschitz estimate without a high-jet hypothesis}
\label{qll-Lipschitz-section}

Let \(\bar\epsilon\) be a representative budget for which the
strict finite two-input comparison holds. It can be chosen so
the base exponent norm \(v\), inherited norms and bad-block
budget have the spare allocations used there; for example
\begin{equation}
 v\le\eta_*/(64e),\qquad
 q(E)\le\eta_*/(128e)
       \quad\hbox{when }\|E\|_{a_0,b}\le\bar\epsilon,
 \qquad 0<\eta_*\le1.
 \label{qll-comparison-budgets}
\end{equation}
The actual linear inherited difference satisfies
\begin{equation}
 q(D)\le16C_TC_R\,\frac{gT}{b}\|D\|_{a_0,b}.
 \label{qll-increment-bound}
\end{equation}
For two quotient classes of norm at most \(\bar\epsilon/4\),
choose covariant near-minimal representatives \(Y\) of \(y\)
and \(D\) of \(x-y\), and set \(X=Y+D\). For small \(\varepsilon>0\),
\begin{align}
 \|Y\|&\le\bar\epsilon/4+\varepsilon,\notag\\
 \|D\|&\le\|x-y\|_{\mathcal Q}+\varepsilon
                    \le\bar\epsilon/2+\varepsilon,\notag\\
 \|X\|&\le3\bar\epsilon/4+2\varepsilon<\bar\epsilon .
 \label{qll-compatible-representatives}
\end{align}
Covariant averaging permits these choices without norm loss.
This construction controls the representative difference;
two unrelated near-minimal representatives would not do so.

For completeness, interpolate the inherited exponent from
\(Q_Y\) to \(Q_Y+zQ_D\). When \(q(D)>0\), use
\begin{equation}
 L_D=\frac{\eta_*}{32e\,q(D)}\ge2.
 \label{qll-interpolation-radius}
\end{equation}
The strict base, inherited and bad-block budgets keep the
Gaussian logarithm family bounded by \(S_1\eta_*\) throughout
this disc. Banach-valued Cauchy gives the finite difference
bound \(64eS_1q(D)\). If \(q(D)=0\), the difference is zero.
Combining with \eqref{qll-increment-bound} yields, after
local limits, Fatou and \(\varepsilon\downarrow0\),
\begin{equation}
 \|\mathcal G_g(x)-\mathcal G_g(y)\|_{\mathcal Q}
 \le C_{\rm quant}\frac{gT}{b}\|x-y\|_{\mathcal Q},
 \qquad C_{\rm quant}=1024eS_1C_TC_R .
 \label{qll-quantum-Lipschitz}
\end{equation}
The pivot and determinant terms cancel because the classical
system is input independent. No fifth-order or other jet
vanishing of \(D\) was used in this logarithm comparison.

If \(B_R\) denotes the established bounded classical transport
constant at this fixed geometry, then
\begin{equation}
 \|\mathcal F_g(x)-\mathcal F_g(y)\|_{\mathcal Q}
 \le\left(B_R+C_{\rm quant}\frac{gT}{b}\right)
                                  \|x-y\|_{\mathcal Q}.
 \label{qll-full-Lipschitz}
\end{equation}
The finite comparison uses the same compatible carrier
assignments on both inputs. Its difference family converges
locally by the construction just proved, so Fatou applies
to that family before the quotient infimum is taken.
This is a small inherited-input dependence of the quantum
remainder; it is not contraction of the whole unextracted map.

\subsection{Varying actual inputs and finite-history compatibility}
\label{qll-varying-section}

\begin{theorem}[Varying-input compatibility]
\label{qll-varying-theorem}
Let finite input families \(E_N\) have exactly invariant
normalized totals, uniform strict budgets and a locally normal
limit \(E\). Use the same input-independent classical system,
fixed windows and exact insertion domains. The actual finite
outputs have local limits by the construction above, and their
limiting normalized class is \(\mathcal F_g([E])\).
The input and output carrier assignments need not be covariant.
\end{theorem}

\begin{proof}
Sections~\ref{qll-Gaussian-section} and
\ref{qll-inherited-section} apply directly to \(E_N\):
they require local input convergence and uniform positive
tails, not the designation \emph{canonical periodization}.
This produces a local output from that particular limiting
carrier family.

To identify its class, average each finite \(E_N\) over its
finite cell-translation group. The norm does not increase.
Its raw finite total is unchanged: normalized translation
invariance implies raw invariance because the zero field
is fixed by translation. Thus the actual finite quantum
output normalized total is also unchanged.

Extract a locally normal subsequence of these covariant
averages, with limit \(E^{\rm cov}\). Their identical finite
normalized totals and interaction compactness imply
\([E^{\rm cov}]=[E]\). The averaged inputs and the canonical
periodizations of \(E^{\rm cov}\) converge locally normally
to the same covariant family. Their input-independent
classical data are the same. Construction continuity therefore
gives identical canonical limiting output families.

Passing the equality between the original and averaged
finite output normalized totals identifies the original
limiting class with \(\mathcal F_g([E])\). This reasoning
applies to any extracted covariant average limit, so no
interaction-dependent replacement of the classical system
is needed. Canonical arrays themselves need not agree.
\end{proof}

By induction, every fixed finite number of actual homogeneous
local steps agrees, at the normalized local-total level, with
the corresponding composition of limiting quotient maps,
provided every intermediate actual finite input remains within
the uniform budgets and analytic insertion domains and the
classical systems and nested regulators are chosen jointly.
At each stage use Theorem~\ref{qll-varying-theorem}, then the
next stage's local continuity. Joint finite lists and countable
fixed parameter systems can be extracted by the earlier
diagonal construction when their hypotheses hold.

All intermediate windows, Gaussian integrations and generated
corrections remain present in this assertion. Each finite
vacuum normalization is tracked separately; the quotient
statement neither discards it from the integral nor proves
convergence of an extensive vacuum energy. This is not an
identification of sequentially cut off integrals with one
one-shot integral.

An infinite quotient relocalization does not by itself provide
small representatives of the actual finite winding remainders.
Thus an infinite quotient bound alone does not supply the
intermediate finite-input hypotheses just stated. If finite
classical reference subtraction is later used, the required
interface is local normal convergence of the actual reference
families and convergence of its extraction coefficients;
equality with a canonical periodized reference is unnecessary
and is not asserted here.

\subsection{Scope}
\label{qll-scope}

The result gives an actual local limit and a continuous
quantum map for the selected homogeneous local sector and
fixed classical system, with the established same-radius
analytic continuation and exact insertion representation.
It is not an assertion about the full Wilson measure or
mixed and large-field histories. It does not prove uniqueness
across arbitrary classical subsequences.

The one-input bounds and the small quantum Lipschitz factor
do not themselves provide an invariant renormalized ball.
A finite block still retains all generated quantum terms;
full-reference extraction, amplitude closure, an actual
coupling trajectory and the remaining field sectors and
observable limits require further work. No infinite
trajectory or continuum construction follows merely from
the finite-history compatibility proved here.
